\documentclass[11pt]{article}

% Prefix for numedquestion's
\newcommand{\questiontype}{Question}


% Use this if your "written" questions are all under one section
% For example, if the homework handout has Section 5: Written Questions
% and all questions are 5.1, 5.2, 5.3, etc. set this to 5
% Use for 0 no prefix. Redefine as needed per-question.
\newcommand{\writtensection}{0}

\usepackage{amsmath, amsfonts, amsthm, amssymb}  % Some math symbols
\usepackage{mathtools}

\usepackage{centernot}
\usepackage{mathtools}

\usepackage{enumitem}

\setlength{\parindent}{0pt}

\begin{document}

\section{Oblivious Transfers}
The beginning of this lectures notes discuss the idea of Oblivious Transfers. Rabin's oblivious transfer protocol is as follows:
\begin{enumerate}
    \item Sender $S$ sends a bit $b$ of information through the OT.
    \item With probability $1/2$, the receiver $R$ receives bit $b$ and with probability $1/2$, it receives a $\#$ symbol which represents loss of information.
\end{enumerate}
The \textbf{1-2 OT} is another protocol that does the following:
\begin{enumerate}
    \item Sender $S$ sends two bits $b_0, b_1$ into the OT channel.
    \item Receiver $R$ submits either $s=0$ or $s=1$ and receives $b_s$ from the OT. 
\end{enumerate}

\subsection{Clear Channels}
Each of the above gadgets imply a clear channel, that is, a channel where information can be passed sucessfully. Firstly, for the 1-2 OT we have that the sender can simply put the same bit for both $b_0$ and $b_1$, or that $S$ and $R$ agree to only use $b_0$ for communication. For Rabin's OT, we have that the sender $S$ can simply send the same information many times, and with overwhelming probability the information is transfered successfully to $R$.

\subsection{1-2 OT $\Longleftrightarrow$ Rabin OT}
The above two gadgets are equivalent, meaning that given either one of the two, the other can be constructed.
\bigskip

\textbf{1-2 OT $\Longrightarrow$ Rabin OT}:\\
Firstly, we will demonstrate that given a $1-2$ OT, we can consruct a Rabin OT easily. We assume that $R$ and $S$ have a clear channel to communicate with one another with and that both parties are honest. The construction is as follows:
\begin{enumerate}
    \item $S$ generates a random value $r$ and flips coin $c$. If $c=0$, it sets $b_0 = b$ and $b_1 = r$ if $c=1$, it sets $b_0=r$ and $b_1=b$.
    \item $R$ guesses either $s=0$ or $s=1$ and receives value $v = b/r$.
    \item $S$ tells $R$ the coin flip result for $c$, indicating whether $s=0/1$ would result in the correct value being received.
\end{enumerate}
Since $R$ has probability $1/2$ of guessing the correct value for $s$, it simulates the Rabit OT properties.
\bigskip

\textbf{Rabin OT $\Longrightarrow$ 1-2 OT}:\\
We demonstrate how to generate a $1-2$ OT given a Rabin OT. For sufficient security parameter $k$, it requires $3k$ rounds of the Rabin OT and is performed as follows:
\begin{enumerate}
    \item $S$ sends $3k$ bits to $R$ which will be represented as $r_1,...,r_{3k}$. The expected number of received bits is $1.5k$ bits for different $r_i$ values and $1.5k\ \#$ symbols.
    \item Following Chernoff bounds, the probability that $R$ receives $\geq 2k$ $r_i$ values is negligible and the probability it receives $\leq k$ $r_i$ values is also negligible, so it must have at least $k$ $r_i$ values sent from $S$.
    \item $R$ selects either $s=0$ or $s=1$ as follows. It selects two disjoint sets of size $k$ of indices for $r_i$ values such that one of the two is composed entirely of $i$ values where it successfully received $r_i$. Notice that there is a overwhelming probability that it received less that $2k$ bits, so it can only construct a single subset of length $k$ such that it knows every $r_i$ value. It sends these two vectors of length $k$ as $(v_0,v_1)$ with its selection $s$ being the $v_s$ vector that knows all corresponding $r_i$ values.
    \item $S$ constructs, using each received vector $\bigoplus_i r_i \oplus b_0$ for $v_0$ and $\bigoplus_i r_i \oplus b_1$ for $v_1$ and sends these values to $R$. Note that the above notation is saying, for each index in $v_s$, we xor all $r_i$ values sent to $R$ and then xor that value with the corresponding $b_s$.
    \item $R$ receives two values, say $v_0, v_1$. Since it knows all the used $r_i$ values for $v_s$, it can calculate $b_s = \bigoplus_i r_i \oplus v_s$ to get the correct bit value.
\end{enumerate}

\section{Implementing 1-2 OT}
The following section introduces constructions for $1-2$ OT.

\subsection{Honest Receiver Construction}
The following construction is easy to break if the receiver is malicious, however, it introduces the main idea of the correct protocol:
\begin{enumerate}
    \item Firstly, the sender $A$ selects a trapdoor permutation $f: \{0,1\}^n \rightarrow \{0,1\}^n$ as well as a random string $r$ to $B$.
    \item $B$ generates its selection $s$ by choosing two $n$-length bitstrings $(y_0,y_1)$. For its selection $s$, it generates $y_s$ by choosing random $n$-length bitstring $x_s$ and computing $y_s = f(x_s)$.
    \item $A$ computes $(b_0 \oplus \langle f^{-1}(y_0), r\rangle, b_1 \oplus \langle f^{-1}(y_1), r\rangle)$. These two values are sent to $B$.
    \item Since an honest $B$ can only have the inverse to $y_0$ or $y_1$, it can only correctly compute $b_0$ or $b_1$.
\end{enumerate}
The issue with this scheme is that it is very easy for $B$ to cheat. They simply need to compute both $y_0$ and $y_1$ by first calculating the input $x_0, x_1$. The following is an extension of this that ensures they cannot do this.
\subsection{Coin Flips Into the Well}
The following protocol produces random bits dictated by Alice, however, Alice cannot see the result of the coinflips. This means that the randomness used by Bob cannot be determined by Bob, however, Alice does not know that randomness he uses. The protocol is as follows:
\begin{enumerate}
    \item Bob selects random bit $R_b = r_{1b}...r_{nb}$.
    \item Alice sends $n$ random bits $R_A =r_{1a}...r_{na}$ to Bob.
    \item Bob's randomness is $R = r_{1b} \oplus r_{1a},..., r_{nb} \oplus r_{na}$.
\end{enumerate}
This makes it very easy for Bob to cheat, however, so the following protocol is instead used in order to ensure that Bob uses the randomness generated by Alice.
\subsubsection{ZK 1-2 OT}
To ensure that the receiver (Bob) cannot cheat, the following protocol is used. Only the coin flips into the well protocol is described as the rest follows from the honest protocol. The bits are generated as follows:
\begin{enumerate}
    \item Bob sends commitments for $r_{b1}...r_{bn}$ as $com(r_{b1}),...,com(r_{bn})$.
    \item Alice sends randomness $r_{a1}...r_{an}$ and sends it to Bob.
    \item Bob retrieves $f,r$ from Alice and returns $y_1,y_2$ generated from random bits $r_{b1} \oplus r_{a1},...,r_{bn} \oplus r_{an}$.
\end{enumerate}
Now, Bob has to prove to Alice that he used her bits in order to generate the randomness according to the protocol. To do this, first consider the following statement:
$$\exists \text{ decommitments to } com(r_{1b})...com(r_{nb}) \text{ s.t.} y_1, y_2 \text{ are generated properly}$$
Consider the witness $decom(r_{1b})...decom(r_{nb})$, that just decommits all the $r_b$ values. The key fact is that this witness can be verified in polynomial time since an algorithm can run to just recompute $y_1,y_2$ given the $r$ values. This means that the above language in \textbf{in NP!} Since this is the case, we can use a ZK proof for a NP complete problem (Ham path, Three Coloring) to confirm that the values $y_1, y_2$ were generated honestly, which ensures that Bob used the coinflips Alice generated.

\section{Secure Computation for N Parties}
We can simulate the secure computation among $n$ honest players for any function $f: \{0,1\}^{n} \rightarrow \{0,1\}$. We use the following theorem:
\bigskip

Any polynomially computatable function $f: \{0,1\}^{n \cdot |x|} \rightarrow \{0,1\}$ can be represented as a polynomial size circuit consisting of only XOR and AND gates.
\bigskip

The above construction involves converting $f$ into a turing machine, the turing machine into a tableu, and the tableu into a circuit.

\subsection{Setup}
The setup of this involves creating shares of the input. This is done among $n$ players $P_1,...,P_n$ by having each player send the other $n-1$ players a random bit $r_{(j,l)}$. Now, each player's share of each input bit $a_i$ is all sent random bits XORed with all received random bits XORed with their actual input $x_{j,i}$.

\subsection{Gate Construction}
The gates follow nearly exactly from the 2-party construction. For XOR gates, the XOR of all shares of inputs $A$ and $B$ creates $C$, in other words:
\begin{align*}
    A \oplus B &= C\\
    (a_1 \oplus a_2 ... \oplus a_n)\oplus(b_1 \oplus ... \oplus b_n) &= (c_1 \oplus ... \oplus c_n)\\
    (a_1 \oplus b_1) \oplus ... \oplus (a_n \oplus b_n) = C
\end{align*}
Therefore, we can just have the players XOR their indivdual shares to produce shares of the output of an XOR gate.
\bigskip

For AND gates, we can decompose the gates into the following:
$$(a_1 \oplus a_2 \oplus ... \oplus a_n) \cdot (b_1 \oplus b_2 \oplus ... \oplus b_n) = \sum_{i=1}^n\sum_{j=1}^n a_ib_j$$
Here we have the same sort of distribution as the two party case. All $a_ib_j$ instancess where $i=j$ can be calculated locally for each player. For every other term, players $p_i,p_j$ can use their pairwise communication channel to compute $a_ib_j$ following the 1-2 OT strategy used in the 2 player model.

\section{Yao's Garbled Circuits}
Garbled circuits involve two parties, a Garbler $G$ and evaluator $E$. The interaction is executed as follows:
\begin{enumerate}
    \item Garbler has randomness $r$ over which they generate garble circuit $c$ to create $\hat{c}$ or $G(c) = \hat{c}$ and input $\hat{x} = G(x)$. They send $\hat{c}$ and $\hat{x}$ to evaluator $E$.
    \item Evaluator runs the function $\hat{c}(\hat{x})$ and the output must be equal to $c(x)$. However, the evaluator must not be able to learn anything about the input other than what the output reveals.
\end{enumerate}
The basics of Yao's strategy for accomplishing this while maining security under the assumption that only one input is given to the user is as follows.
\begin{enumerate}
    \item First we assume the entire circuit is composed of gates that take in 2 bits and output a single bit. We also assume the circuit's final output is a single bit.
    \item We define each wire as being a $0/1$ by pair of keys $k_0, k_1$, each key corresponding to the respective bit.
    \item These keys are represented as a table, and each gate with inputs $a_0,b_1$ defines its output as encryptions of $k_b$ using the keys from its input wires.
    \item To evaluate a gate, the evaluator decrypts the corresponding entry in the table that is encrypted by the two input keys they got from the previous wire (The table must be permuted randomly which is explained in the next set of notes)
    \item The final wire's output is either $k_0$ or $k_1$ which the garbler can decrypt as $0/1$
\end{enumerate}



\end{document}